% !TEX spellcheck = en_US
% !TEX spellcheck = LaTeX
\documentclass[letterpaper,10pt,english]{article}
\usepackage{%
amsfonts,%
amsmath,%
amssymb,%
amsthm,%
babel,%
bbm,%
%biblatex,%
caption,%
centernot,%
color,%
enumerate,%
%enumitem,%
epsfig,%
epstopdf,%
etex,%
fancybox,%
framed,%
fullpage,%
%geometry,%
graphicx,%
hyperref,%
latexsym,%
mathptmx,%
mathtools,%
multicol,%
paralist,%
pgf,%
pgfplots,%
pgfplotstable,%
pgfpages,%
proof,%
psfrag,%
%subfigure,%
tcolorbox,%
tfrupee,%
tikz,%
times,%
%ulem,%
url,%
xcolor,%
mathpazo,%
}

\definecolor{shadecolor}{gray}{.95}%{rgb}{1,0,0}
\usepackage[margin=1in,top=0.75in]{geometry}
\usepackage[mathscr]{eucal}
\usepgflibrary{shapes}
\usepgfplotslibrary{fillbetween}
\usetikzlibrary{%
arrows,%
backgrounds,%
chains,%
decorations.pathmorphing,% /pgf/decoration/random steps | erste Graphik
decorations.text,% 
matrix,%
positioning,% wg. " of "
fit,%
patterns,%
petri,%
plotmarks,%
scopes,%
shadows,%
shapes.misc,% wg. rounded rectangle
shapes.arrows,%
shapes.callouts,%
shapes%
}

%\pgfplotsset{compat=newest} %<------ Here
\pgfplotsset{compat=1.11} %<------ Or use this one

\theoremstyle{plain}
\newtheorem{thm}{Theorem}[section]
\newtheorem{lem}[thm]{Lemma}
\newtheorem{prop}[thm]{Proposition}
\newtheorem{cor}[thm]{Corollary}
\newtheorem{clm}[thm]{Claim}

\theoremstyle{definition}
\newtheorem{defn}[thm]{Definition}
\newtheorem{conj}[thm]{Conjecture}
\newtheorem{exmp}[thm]{Example}
\newtheorem{axiom}[thm]{Axiom}
\newtheorem{assum}[thm]{Assumption}
\newtheorem{exerc}[thm]{Exercise}
\newtheorem*{defin}{Definition}

\theoremstyle{remark}
\newtheorem{rem}{Remark}
\newtheorem{note}{Note}

\newcommand{\iid}{\emph{i.i.d.}\ }
\newcommand{\Cov}{\operatorname{Cov}}
%\newcommand{\det}{\operatorname{det}}
%\newcommand{\Real}{\mathbb{R}}
\newcommand{\tr}{\operatorname{tr}}
\newcommand{\Var}{\operatorname{Var}}
\newcommand{\cov}{\operatorname{cov}}

%\renewcommand{\proof}[1]{\begin{proof}#1\end{proof}}
\newcommand{\EQ}[1]{\begin{equation*}#1\end{equation*}}
\newcommand{\EQN}[1]{\begin{equation}#1\end{equation}}
\newcommand{\eq}[1]{\begin{align*}#1\end{align*}}
\newcommand{\meq}[2]{\begin{xalignat*}{#1}#2\end{xalignat*}}
\newcommand{\norm}[1]{\left\lVert#1\right\rVert}
\newcommand{\inner}[1]{\left\langle #1\right\rangle}
\newcommand{\abs}[1]{\left\lvert#1\right\rvert}
\newcommand{\expect}[1]{\mathbb{E}\left[{#1}\right]}
\newcommand{\prob}[1]{\mathbb{P}\left[{#1}\right]}
\newcommand{\given}{\; \big\vert \;} 
\newcommand{\set}[1]{\left\{#1\right\}} 
\newcommand{\indicator}[1]{1_{\set{#1}}} 
\newcommand{\Ind}[1]{\mathbbm{1}_{#1}} 
\newcommand{\SetIn}[1]{\mathbbm{1}_{\set{#1}}} 
\newcommand{\red}[1]{\textcolor{red}{#1}} 
\newcommand{\floor}[1]{\left\lfloor#1\right\rfloor}
\newcommand{\ceil}[1]{\left\lceil#1\right\rceil}


\newcommand{\C}{\mathbb{C}}
\newcommand{\D}{\mathbb{D}}
\newcommand{\E}{\mathbb{E}}
\newcommand{\F}{\mathbb{F}}
\newcommand{\N}{\mathbb{N}}
\renewcommand{\P}{\mathbb{P}}
\newcommand{\Q}{\mathbb{Q}}
\newcommand{\R}{\mathbb{R}}
\newcommand{\Z}{\mathbb{Z}}

\newcommand{\bA}{\mathbf{A}}
\newcommand{\bI}{\mathbf{I}}
\newcommand{\bU}{\mathbf{1}}
\newcommand{\bx}{\mathbf{x}}
\newcommand{\bX}{\mathbf{X}}
\newcommand{\bZ}{\mathbf{Z}}


\newcommand{\cA}{\mathcal{A}}
\newcommand{\cB}{\mathcal{B}}
\newcommand{\cC}{\mathcal{C}}
\newcommand{\cF}{\mathcal{F}}
\newcommand{\cG}{\mathcal{G}}
\newcommand{\cM}{\mathcal{M}}
\newcommand{\cN}{\mathcal{N}}
\newcommand{\cP}{\mathcal{P}}
\newcommand{\cT}{\mathcal{T}}
\newcommand{\cX}{\mathcal{X}}
\newcommand{\cY}{\mathcal{Y}}

\newcommand{\sA}{\mathscr{A}}
\newcommand{\sB}{\mathscr{B}}
\newcommand{\sC}{\mathscr{C}}
\newcommand{\sE}{\mathscr{E}}
\newcommand{\sF}{\mathscr{F}}
\newcommand{\sG}{\mathscr{G}}
\newcommand{\sH}{\mathscr{H}}
\newcommand{\sL}{\mathscr{L}}
\newcommand{\sS}{\mathscr{S}}
\newcommand{\sT}{\mathscr{T}}
\newcommand{\sX}{\mathscr{X}}
\newcommand{\sY}{\mathscr{Y}}
\newcommand{\sZ}{\mathscr{Z}}

\newcommand{\tS}{\tilde{S}}
% Debug
\newcommand{\todo}[1]{\begin{color}{blue}{{\bf~[TODO:~#1]}}\end{color}}

% a few handy macros

\renewcommand{\le}{\leqslant}
\renewcommand{\ge}{\geqslant}
\newcommand\matlab{{\sc matlab}}
\newcommand{\goto}{\rightarrow}
\newcommand{\bigo}{{\mathcal O}}
%\newcommand{\half}{\frac{1}{2}}
%\newcommand\implies{\quad\Longrightarrow\quad}
\newcommand\reals{{{\rm l} \kern -.15em {\rm R} }}
\newcommand\complex{{\raisebox{.043ex}{\rule{0.07em}{1.56ex}} \hskip -.35em {\rm C}}}


% macros for matrices/vectors:

% matrix environment for vectors or matrices where elements are centered
\newenvironment{mat}{\left[\begin{array}{ccccccccccccccc}}{\end{array}\right]}
\newcommand\bcm{\begin{mat}}
\newcommand\ecm{\end{mat}}

% matrix environment for vectors or matrices where elements are right justifvied
\newenvironment{rmat}{\left[\begin{array}{rrrrrrrrrrrrr}}{\end{array}\right]}
\newcommand\brm{\begin{rmat}}
\newcommand\erm{\end{rmat}}

% for left brace and a set of choices
%\newenvironment{choices}{\left\{ \begin{array}{ll}}{\end{array}\right.}
\newcommand\when{&\text{if~}}
\newcommand\otherwise{&\text{otherwise}}
% sample usage:
%\delta_{ij} = \begin{choices} 1 \when i=j, \\ 0 \otherwise \end{choices}


% for labeling and referencing equations:
\newcommand{\eql}{\begin{equation}\label}
\newcommand{\eqn}[1]{(\ref{#1})}
% can then do
%\eql{eqnlabel}
%...
%\end{equation}
% and refer to it as equation \eqn{eqnlabel}.
% some useful macros for finite difference methods:
\newcommand\unp{U^{n+1}}
\newcommand\unm{U^{n-1}}
% for chemical reactions:
\newcommand{\react}[1]{\stackrel{K_{#1}}{\rightarrow}}
\newcommand{\reactb}[2]{\stackrel{K_{#1}}{~\stackrel{\rightleftharpoons}
 {\scriptstyle K_{#2}}}~}
\makeatletter
\def\th@plain{%
\thm@notefont{}% same as heading font
\itshape % body font
}
\def\th@definition{%
\thm@notefont{}% same as heading font
\normalfont % body font
}
\makeatother
\date{}


\title{Lecture-13: Conditional expectation}
\author{}

\begin{document}
\maketitle


\section{Conditional expectation conditioned on a random  vector}

Let $X: \Omega \to \R^n$ be a random vector defined on a probability space $(\Omega, \sF, P)$. 
Recall that the projection $\pi_i: \R^n \to \R$ of $n$-dimensional vector $X$ on its $i$th component gives $X_i = \pi_i(X)$, 
and it is a Borel measurable function. 
It follows that $X_i = \pi_i\circ X$ is a random variable for each $i \in [n]$. 
Further the event space generated by random vectors $X$ is given by 
\EQ{
\sigma(X) = \sigma(X_1, \dots, X_n).
}
%\begin{defn}[Event space generated by a random vector] 
%Let $X: \Omega \to \R^n$ be a random vector, then the smallest event space generated by the events of the form 
%\EQ{
%A_X(x) 
%\triangleq X^{-1}(-\infty, x_1]\times \dots \times (-\infty, x_n] 
%= \cap_{i=1}^nX_i^{-1}(-\infty, x_i],\quad x \in \R^n, 
%} 
%is called the \textbf{event space generated} by this random vector $X$, and denoted by $\sigma(X)$. 
%\end{defn}
%\begin{shaded*}
%\begin{exmp}[Indicator functions] 
%Let $A, B \in \sF$ be events, then $X = (\Ind{A}, \Ind{B})$ is a random vector. 
%For $x \in \R^2$, we have  
%\EQ{
%%X^{-1}(-\infty, x_1]\times(-\infty, x_2] = X_1^{-1}(-\infty, x_1]\cap X_2^{-1}(-\infty, x_2]
%A(x) = \begin{cases}
%\Omega, &\min\set{x_1, x_2} \ge 1,\\
%A^c, & x_1 \in [0, 1),x_2 \ge 1,\\
%B^c, & x_1 \ge 1,x_2 \in [0,1),\\
%A^c\cap B^c, & x_1, x_2 \in [0, 1),\\
%\emptyset, &\min\set{x_1, x_2}< 0.
%\end{cases}
%}
%This implies that the smallest event space generated by this random vector is 
%\EQ{
%\sigma(X) = \set{\emptyset, A, A^c, B, B^c, A\cup B, A\cup B^c, A \cap B,  A\cap B^c, A^c\cup B, A^c\cup B^c, A^c \cap B,  A^c\cap B^c,\Omega} = \sigma(A, B).
%} 
%\end{exmp}
%\end{shaded*}
\begin{rem}
Let $(A_i \in \sF:  i \in [n])$ be an $n$-length sequence of events, 
then $X = (\Ind{A_i}: i \in [n])$ is a random vector, and the smallest event space generated by this random vector is $\sigma(X) = \sigma(A_1, \dots, A_n)$. 
\end{rem}
\begin{defn}%[Conditional expectation given a random vector]
Consider a random variable $X:\Omega \to\R$ such that $\E\abs{X} < \infty$ and a random vector $Y:\Omega \to \R^n$ for some $n\in \N$, 
defined on the same probability space $(\Omega, \sF, P)$. 
The \textbf{conditional expectation of random variable $X$ given the random vector $Y$} is defined as $\E[X\mid Y] \triangleq \E[X\mid\sigma(Y)]$. 
\end{defn}
\begin{lem} 
For a random sequence $X: \Omega \to \R^\N$ defined on the same probability space $(\Omega, \sF, P)$, 
we have 
\EQ{
\sigma(X_1, \dots, X_n) \subseteq \sigma(X_1, \dots, X_{n+1}),\quad n\in \N.
}
\end{lem}
\begin{proof}
For any $x \in \R^{n+1}$, 
any generating event for collection $\sigma(X_1, \dots, X_n)$ is of the form $\cap_{i=1}^nX_i^{-1}(-\infty, x_i] = \cap_{i=1}^nX_i^{-1}(-\infty, x_i]\cap X_{n+1}^{-1}(\R)$, a generating event for collection $\sigma(X_1, \dots, X_{n+1})$.  
\end{proof}


%\begin{enumerate}[(i)]
%\item The \textbf{conditional distribution} of the random variable $X$ given a non-trivial event $B$ generated by the random vector $Y$ is 
%\EQ{
%F_{X \mid B}(x) = \frac{P(\set{X \le x}\cap B)}{P(B)}. 
%}
%We can define $F_{X\mid B} = 0$ for trivial events $B \in \sigma(Y)$ such that $P(B) = 0$. 
%
%\item The \textbf{conditional expectation} of random variable $X$ given any event $B$ generated by random vector $Y$ is given by 
%\EQ{
%\E[X \mid B] = \int_{x \in \R}xdF_{X\mid B}.
%}
%%We can define $\E[X\mid B] = 0$ for trivial events $B \in \sigma(Y)$ such that $P(B) = 0$. 
%
%%\item  The \textbf{conditional distribution} of the random variable $X$ given the random vector $Y$, 
%%is a function of the random vector $Y$, and hence a random variable. 
%%In particular, for each event $\set{Y\le y}=\cap_{i=1}^n\set{Y_i \le y_i}$
%%\EQ{
%%F_{X\mid Y}\SetIn{Y \le y} = F_{X \mid \set{Y \le y}}\SetIn{Y \le y} = \frac{F_{X,Y}(\cdot,y)}{F_Y(y)}\SetIn{Y \le y}~\text{ for all }y \in \R^n.
%%}
%%That is, this conditional distribution is the collection $(F_{X \mid B}: B \in \sigma(Y))$ of conditional distributions for each event $B$ generated by the random vector $Y$. 
%%The probability of the event $B$ and $F_{X \mid B}$ is $P(B)$. 
%%
%%\item The \textbf{conditional expectation} of random variable $X$ given the random vector $Y$, 
%%is a function of the random vector $Y$, and hence a random variable. 
%%In particular, for each event $\set{Y \le y}=\cap_{i=1}^n\set{Y_i \le y_i}$, we have the event 
%%\EQ{
%%\E[X\mid Y]\SetIn{Y \le y} = \E[X \mid \set{Y \le y}]\SetIn{Y \le y} = \SetIn{Y \le y}\int_{x \in \R}x dF_{X\mid \set{Y \le y}}(x),\text{ for all }y \in \R^n.
%%}
%%That is, this conditional expectation is the collection of events $ (\E[X|B]: B \in \sigma(Y))$ corresponding to the  conditional expectation given the events generated by the random vector $Y$. 
%%The probability of the event $B$ and $\E[X\mid B]$ is $P(B)$. 
%\item The conditional expectation $\E[X\mid Y]: \Omega \to \R$ is Borel measurable function of random vector $Y$.  
%\item The mean of random variable $\E[X\mid Y]\SetIn{Y \le y}$ for all $y \in \R^n$ is given by $\E[X\SetIn{Y \le y}]$.  
%%\EQ{
%%\E[\E[X\mid Y]\SetIn{Y \le y}] 
%%= \int_{t \in \R^n}dF_Y(t)\SetIn{Y \le y}\int_{x \in \R}x dF_{X\mid \set{Y \le y}}(x) 
%%= \int_{(x,t) \in \R^{n+1}}x \SetIn{t \le y}dF_{X,Y}(x,t)
%%= \E[X\SetIn{Y \le y}].
%%}
%Similarly, for all events $B \in \sigma(Y)$, we have 
%\EQ{
%\E[\E[X\mid Y]\Ind{B}] = \E[X\Ind{B}].
%}
%\end{enumerate}

%\begin{rem} 
%There are three main takeaways from this definition. 
%For a random vector $Y:\Omega\to \R^n$ defined on the probability space $(\Omega, \sF, P)$, 
%the event space generated by $Y$ is $\sigma(Y)$.  
%Let $X: \Omega \to \R$ be a random variable on the same probability space.  
%\begin{enumerate}
%\item 
%The conditional expectation $\E[X|Y] = \E[X|\sigma(Y)]$ and is $\sigma(Y)$ measurable. 
%That is, $\E[X|Y]: Y \mapsto \R$ is a Borel measurable function of $Y$. 
%In particular, this implies that $\E[X|Y]$ is a random variable, 
%that takes value $\E[X|E_y]$ when $\omega \in E_y$. 
%When $Y$ is simple, $\E[X|Y]$ is a simple random variable with PMF $P_Y$. 
%When $Y$ is continuous, $\E[X|Y]$ is a continuous random variable with density $f_Y$. 
%
%\item Expectation is averaging. 
%Conditional expectation is averaging over event spaces. 
%We can observe that the coarsest averaging is $\E[X|\set{\emptyset, \Omega}] = \E X$ 
%and the finest averaging is $\E[X|\sigma(X)] = X$. 
%Further, $\E[X|\sigma(Y)]$ is averaging of $X$ over events generated by $Y$.   
%If we take any event $A \in \sigma(Y)$ generated by $Y$, 
%then the conditional expectation of $X$ given $Y$ is fine enough to find the avergaging of $X$ when this event occurs.   
%That is, $\E[X\Ind{A}] = \E[\E[X|Y]\Ind{A}]$.
%
%\item If $X \in L^1$, then the conditional expectation $\E[X|Y] \in L^1$. 
%\end{enumerate}
%\end{rem}
%%%%%%%%%%%%%%%%%%%%%%%%%%%%%%%%%%%%%%%%%%%%%%%%%%%%%%%%%%%%
\section{Properties of Conditional Expectation}
%%%%%%%%%%%%%%%%%%%%%%%%%%%%%%%%%%%%%%%%%%%%%%%%%%%%%%%%%%%%

\begin{prop} 
Let $X,Y$ be random variables on the probability space $(\Omega, \sF, P)$ such that $\E\abs{X} , \E\abs{Y} < \infty$. 
Let $\sG$ and $\sH$ be event spaces such that $\sG, \sH \subset \sF$. 
Then the following statements are true. 
\begin{enumerate}
\item \textbf{Identity:} If $X$ is $\sG$-measurable and $\E\abs{X} < \infty$, then $X = \E[X\mid\sG]$ a.s. 
In particular, $c = \E[c\mid \sG]$, for any constant $c \in  \R$. 
\item \textbf{Linearity:} $\E[(\alpha X + \beta Y) \mid \sG] = \alpha \E[X\mid\sG] + \beta E[Y\mid \sG]$ a.s.
\item \textbf{Monotonicity:} If $X \le Y$ a.s., then $\E[X\mid\sG] \le E[Y\mid \sG]$ a.s.
\item \textbf{Conditional Jensen's inequality:} If $\psi : \R \to  \R$ is convex and $\E\abs{\psi(X)} <\infty$, 
then $\psi(\E[X\mid\sG]) \le \E[\psi(X) \mid \sG]$ a.s. 
%\item \textbf{$L^p$-nonexpansivity:} If $X \in  L^p$, for $p \in  [1, \infty]$, then $\E[X\mid\sG] \in  L^p$
%and $\norm{\E[X\mid\sG]}_p \le \norm{X}_p$.
%In particular, $\E[\abs{X} \mid \sG] \ge \abs{\E[X\mid\sG]}$ a.s.

\item \textbf{Pulling out what's known:} If $Y$ is $\sG$-measurable and $\E\abs{XY} <\infty$, then
$\E[XY\mid \sG] = Y\E[X\mid\sG]$ a.s.
\item \textbf{Tower property:} If $\sH \subseteq \sG$, then
$\E[\E[X\mid\sG]\mid \sH] = \E[X\mid \sH]$ a.s.
\item \textbf{Irrelevance of independent information:} If $\sH$ is independent of $\sigma (\sG, \sigma (X))$
then
$\E[X|\sigma (\sG, \sH)] = \E[X\mid\sG] \text{ a.s. }$
In particular, if $X$ is independent of $\sH$, then $\E[X\mid \sH] = \E[X]$ a.s.
\item \textbf{$L^2$-projection:} If $\E\abs{X}^2 < \infty$, then $\zeta^\ast \triangleq \E[X\mid\sG]$ minimizes $\E[(X - \zeta)^2]$
over all $\sG$-measurable random variables $\zeta$ such that $\E\abs{\zeta}^2 < \infty$. 
\end{enumerate}
\end{prop}
\begin{proof}
%Let $X,Y\in L^1$ be random variables on the probability space $(\Omega, \sF, P)$. 
%Let $\sG$ and $\sH$ be event spaces such that $\sG, \sH \subseteq \sF$. 
In all the properties below, we have to show that conditional expectation given an event space equals another random variable almost surely. 
To this end, we will show that the right hand side satisfies the three properties of the conditional expectation random variable, 
and hence is the conditional expectation from almost sure uniqueness.  
\begin{enumerate}
\item \textbf{Identity:} It follows from the definition that $X$ satisfies all three conditions for conditional expectation.  
The event space generated by any constant function is the trivial event space $\set{\emptyset, \Omega} \subseteq \sG$ for any event space. 
Hence, $\E[c\mid \sG] = c$.  
\item \textbf{Linearity:} 
Let $Z \triangleq \E[(\alpha X + \beta Y) \mid \sG]$ and $Z_1 \triangleq \E[X\mid\sG]$ and $Z_2\triangleq\E[Y\mid\sG]$. 
We have to show that $Z = \alpha Z_1+ \beta Z_2$ almost surely. 
\begin{enumerate}[(i)]
\item From the definition of conditional expectation, we have $\E[X\mid\sG], \E[Y\mid\sG]$ are $\sG$-measurable. 
In addition, the linear map $g: \R^2\to\R$ defined by $g(x) \triangleq \alpha x_1 + \beta x_2$ is Borel measurable. 
Therefore, the linear combination $\alpha Z_1 + \beta Z_2$ is $\sG$-measurable. 
\item For any $G \in \sG$, from the linearity of expectation and definition of conditional expectation, we have 
\EQ{
\E[(\alpha Z_1 + \beta Z_2)\Ind{G}] 
= \alpha\E[ \E[X\mid \sG]\Ind{G}] + \beta\E[ \E[Y\mid \sG]\Ind{G}] 
= \E[(\alpha X + \beta Y)\Ind{G}].
} 
\item From the definition of conditional expectation, we have $\E\abs{Z_1}, \E\abs{Z_2}$ are finite. 
Therefore, we have 
\EQ{
\E\abs{\alpha Z_1 + \beta Z_2} \le \abs{\alpha}\E\abs{Z_1} + \abs{\beta}\E\abs{Z_2} < \infty.
}
\end{enumerate}
This implies that $\alpha Z_1 + \beta Z_2$ satisfies three properties of conditional expectation $\E[(\alpha X +\beta Y)\mid\sG]$. 
From the almost sure uniqueness of conditional expectation, we have $Z = \alpha Z_1 + \beta Z_2$ almost surely. 

\item \textbf{Monotonicity:} Let $\epsilon > 0$ and define $A_\epsilon \triangleq \set{\E[X\mid\sG] - \E[Y\mid \sG] >\epsilon} \in  \sG$.  
Then from the definition of conditional expectation, we have 
\EQ{
0 \le \epsilon P(A_\epsilon) 
= \E[\epsilon\Ind{A_\epsilon}] 
< \E[(\E[X\mid \sG]-\E[Y\mid\sG])\Ind{A_\epsilon}] 
= \E[(X-Y)\Ind{A_\epsilon}] 
\le 0.
}
Thus, we obtain  that $P(A_\epsilon) = 0$ for all $\epsilon > 0$. 
Taking limit $\epsilon \downarrow 0$, we get 
\EQ{
0 =\lim_{\epsilon \downarrow 0}P(A_\epsilon) = P(\lim_\epsilon A_\epsilon) = P(A_0).
} 

\item \textbf{Conditional Jensen's inequality:}
We will use the fact that a convex function can always be represented by the supremum of a family of affine functions. 
Accordingly, we will assume for a convex function $\psi:\R\to\R$, 
we have linear functions $\phi_i:\R\to\R$ and constants $c_i\in\R$ for all $i \in I$ such that $\psi = \sup_{i\in I}(\phi_i+c_i)$. 

For each $i\in I$, we have $\phi_i(\E[X\mid\sG])+c_i = \E[\phi_i(X)\mid\sG] +c_i\le \E[\psi(X)\mid\sG]$ from the linearity and monotonicity of conditional expectation. 
It follows that
\EQ{
\psi(\E[X\mid\sG]) = \sup_{i \in I}(\phi_i(\E[X\mid\sG]) + c_i) \le \E[\psi(X)\mid\sG].
}

%Use the result of Lemma~\ref{} which states that $\psi(x) = \sup_{n\in \N}(a_n + b_nx)$, where $(a_n \in \R: n \in \N)$ and $(b_n: n\in \N)$ are sequences of real numbers. 

%\item \textbf{$L^p$-nonexpansivity:} For $p \in  [1, \infty)$, apply conditional Jensen's inequality with $\psi(x) = \abs{x}^p$. The case $p = \infty$ follows directly.

\item \textbf{Pulling out what's known:} 
From the almost sure uniqueness of conditional expectation,  
it suffices to show that $Y\E[X\mid\sG]$ satisfies following three properties of the conditional expectation $\E[XY\mid\sG]$: 
\begin{enumerate}[(i)]
\item the random variable $Y\E[X\mid\sG]$ is $\sG$-measurable,
\item $\E[XY\Ind{G}] = \E[Y\E[X\mid\sG]\Ind{G}]$ for any event $G\in\sG$, and 
\item $\E\abs{Y\E[X\mid\sG]}$ is finite.
\end{enumerate}
Part~(i) is true since $Y$ is given to be $\sG$-measurable, 
$\E[X\mid\sG]$ is $\sG$-measurable by the definition of conditional expectation, 
and product function is Borel measurable. 

It suffices to show part~(ii) and~(iii) for simple $\sG$-measurable random variables $Y=\sum_{y\in\sY}y\Ind{E_y}$ where $E_y = Y^{-1}\set{y} \in \sG$. %, where $\E\abs{XY}$ is finite. 
For any $G\in\sG$, we have $G\cap E_y\in\sG$ for all $y \in \sY$ and $\cup_{y\in\sY}(G\cap E_y) = G$. 

Part~(ii) follow from the linearity of expectation and definition of conditional expectation $\E[X\mid\sG]$, 
such that  
\EQ{
\E[Y\E[X\mid\sG]\Ind{G}] 
= \sum_{y\in\sY}y\E[\Ind{G\cap E_y}\E[X\mid\sG]]
= \sum_{y\in\sY}y\E[X\Ind{G\cap E_y}]
= \E[X\sum_{y\in\sY}y\Ind{G\cap E_y}]
= \E[XY\Ind{G}].
}

Part~(iii) follows from the fact that $\E\abs{XY}$ is finite and the conditional Jensen's inequality applied to convex function $\abs{\cdot} : \R \to \R_+$ to get $\abs{\E[X\mid\sG]} \le \E[\abs{X}\mid\sG]$. 
Therefore, 
\EQ{
\E[\abs{Y}\abs{\E[X\mid\sG]}] 
= \sum_{y\in\sY}\abs{y}\E[\abs{\E[X\mid\sG]}\Ind{E_y}]
\le \sum_{y \in \sY}\abs{y}\E[\abs{X}\Ind{E_y}] 
= \E\abs{XY}.  
} 

%it suffices to show that $\E[ZX] = \E[Z\E[X\mid\sG]]$ for any simple $\sG$-measurable random variable $Z$ with $\E\abs{ZX} < \infty$, 
%from which the previous statement follows for $Z = Y\Ind{A}$. 
%
%%Part~(iii) is true since $\E\abs{Y\E[X\mid\sG]} \le $ and $\E\abs{XY} < \infty$, 
%%then we need to show that $\E[XY\Ind{A}] = \E[Y\E[X\mid\sG]\Ind{A}]$, for all $A \in  \sG$. 
%
%
%Let $Z = \sum_{k=1}^n\alpha_k\Ind{A_k}$ for $(A_1, \dots, A_n) \subset \sG$, then the result is a consequence of the definition of conditional expectation and linearity. 
%
%%Let us assume that both $Z$ and $X$ are nonnegative and $ZX \in  L^1$. 
%%In that case we can find a nondecreasing sequence $(Z_n \ge 0: n\in \N)$ of non-negative simple random variables with $Z_n \uparrow Z$. 
%%Then $Z_nX \in  L^1$ for all $n \in  \N$ and the monotone convergence theorem implies that 
%%\EQ{
%%\E[ZX] = \lim_n\E[Z_nX] = \lim_n\E[Z_n\E[X\mid\sG]] = \E[Z\E[X\mid\sG]].
%%}
%%Our next task is to relax the assumption $X \in  L^1_+$ to the original one $X \in  L^1$. 
%%In that case, the $L^p$-nonexpansivity for p = 1 implies that $\abs{\E[X\mid\sG]} \le \E[\abs{X} \mid \sG]$ a.s.,
%%and so
%%$\abs{Z_n\E[X\mid\sG]} \le Z_n\E[\abs{X} \mid \sG] \le ZE[\abs{X} \mid \sG]$.
%%We know from the previous case that
%%$\E[Z\E[\abs{X} \mid \sG]] = \E[Z \abs{X}]$, so that $Z\E[\abs{X} \mid \sG] \in  L^1$.
%%We can, therefore, use the dominated convergence theorem to conclude that 
%%\EQ{
%%\E[Z\E[X\mid\sG]] = \lim_n\E[Z_n\E[X\mid\sG]] = \lim_n\E[Z_nX] = \E[ZX].
%%}
%%Finally, the case of a general $Z$ follows by linearity.
\item \textbf{Tower property:} 
From the almost sure uniqueness of conditional expectation, 
it suffices to show that 
\begin{enumerate}[(i)]
\item $\E[X\mid\sH]$ is $\sH$-measurable, 
\item $\E[\E[X\mid\sH]\Ind{H}] = \E[\E[X\mid\sG]\Ind{H}]$ for all $H\in\sH$, and 
\item $\E[\abs{\E[X\mid\sH]}]$ is finite.  
\end{enumerate}
Part~(i) follows from the definition of conditional expectation, 
which implies that $\E[X\mid \sH]$ is $\sH$ measurable. 

Part~(ii) follows from the fact that $\sH \subseteq \sG$, 
and hence any $H \in \sH$ belongs to $\sG$. 
Therefore, from the definition of conditional expectation, we have 
\EQ{
\E[\E[X\mid \sG]\Ind{H}] = \E[X\Ind{H}] = \E[\E[X\mid \sH]\Ind{H}]. 
}

Part~(iii) follows from the conditional Jensen's inequality applied to convex function $\abs{\cdot}:\R\to\R_+$ to get $\abs{\E[X\mid\sH]} \le \E[\abs{X}\mid\sH]$, 
and hence $\E\abs{\E[X\mid\sH]} \le \E\abs{X} < \infty$. 

\item \textbf{Irrelevance of independent information:} 
From the almost sure uniqueness of conditional expectation, 
it suffices to show that 
\begin{enumerate}[(i)]
\item $\E[X\mid\sG]$ is $\sigma(\sG,\sH)$-measurable, 
\item $\E[\E[X\mid\sG]\Ind{G}] = \E[X\Ind{G}]$ for all $G\in\sigma(\sG,\sH)$, and 
\item $\E\abs{\E[X\mid\sG]}$ is finite.  
\end{enumerate}

Part~(i) follows from the definition of conditional expectation and the definition of $\sigma(\sG,\sH)$. 
Since $\E[X\mid\sG]$ is $\sG$-measurable, it is $\sigma(\sG,\sH)$ measurable. 


%We assume $X \ge 0$ and show that
%\EQ{
%\E[X\Ind{A}] = \E[\E[X\mid\sG]\Ind{A}], \text{ a.s. for all }A \in  \sigma (\sG, \sH). 
%}
%%Let $\sL$ be the collection of all $A \in  \sigma (\sG, \sH)$ such that (10.3) holds.
%%It is straightforward that $\sL$ is a $\lambda$-system, so it will be enough to
%%establish (10.3) for some $\pi$-system that generates $\sigma (\sG, \sH)$. 
%%One possibility is $P = \set{G \cap H : G \in  \sG, H \in  \sH}$, and for $G \cap H \in  P$ we use
%%independence of $\Ind{H}$ and $\E[X\mid\sG]\Ind{G}$, as well as the independence of
%%$\Ind{H}$ and $X\Ind{G}$ to get 
Part~(ii) follows from the fact that it suffices to show for events $A = G \cap H \in \sigma(\sG,\sH)$ 
where $G \in \sG$ and $H \in \sH$. 
In this case,  
\EQ{
\E[\E[X\mid\sG]\Ind{G\cap H}] = \E[\E[X\mid\sG]\Ind{G}\Ind{H}] 
= \E[\E[X\mid\sG]\Ind{G}]E[\Ind{H}]
= \E[X\Ind{G}]\E[\Ind{H}] = \E[X\Ind{G\cap H}]. 
}

Part~(iii) follows from the conditional Jensen's inequality applied to convex function $\abs{\cdot}:\R\to\R_+$ to get $\abs{\E[X\mid\sG]} \le \E[\abs{X}\mid\sG]$. 
This implies that $\E\abs{\E[X\mid\sG]} \le \E\abs{X} < \infty$. 

\item \textbf{$L^2$-projection:} 
We define $L^2(\sG) \triangleq \set{\zeta \text{ a } \sG \text{ measurable random variable}: 
\E\zeta^2 < \infty%,\zeta^{-1}(-\infty, x] \in \sG \text{ for all }x\in \R
}$. 
From the conditional Jensen's inequality applied to convex function $(\cdot)^2:\R \to \R_+$, 
we get that $(\E[X\mid\sG])^2 \le \E[X^2\mid\sG]$ a.s.. 
Since $X\in L^2$, it follows that $X^2\in L^1$ and hence $\E[X\mid\sG] \in L^2$. 
It follows that $\zeta^\ast \triangleq \E[X\mid\sG] \in L^2(\sG)$ from the definition of conditional expectation. 
 
We first show that $X - \zeta^\ast$ is uncorrelated with all $\zeta \in L^2(\sG)$. 
Towards this end, we let $\zeta\in L^2(\sG)$ and observe that 
\EQ{
\E[(X - \zeta^\ast)\zeta] 
= \E[\zeta X] - \E[\zeta \E[X\mid\sG]] 
= \E[\zeta X] - \E[\E[\zeta X\mid \sG]] = 0.
}
The above equality follows from the linearity of expectation, 
the $\sG$-measurability of $\zeta$,  
and the definition of conditional expectation. 
Since $\zeta^\ast\in L^2(\sG)$, we have $(\zeta-\zeta^\ast)\in L^2(\sG)$. 
Therefore, $\E[(X - \zeta^\ast)(\zeta-\zeta^\ast)] = 0$.  
 
For any $\zeta\in L^2(\sG)$, 
we can write from the linearity of expectation 
\EQ{
\E(X-\zeta)^2 
= \E(X-\zeta^\ast)^2 + \E(\zeta-\zeta^\ast)^2-2\E(X-\zeta^\ast)(\zeta-\zeta^\ast)
\ge \E(X-\zeta^\ast)^2.
}
%It is enough to show that 
%This implies that $\E(X-\zeta)^2 \ge \E(X-\zeta^\ast)^2$ for all $\sG$ measurable random variables $\zeta$ that have finite second moment. 


\end{enumerate}
\end{proof}


%\begin{shaded*}
%\begin{exmp}
%If $X^{-1}(-\infty, x]\in \sG$ for all $x \in \R$, then $\E[X\mid \sG] = X$. \\
%When, $X$ is given to be a random variable with respect to $\sG$, this follows directly from the definition.
%\end{exmp}
%\end{shaded*}
%
%\begin{shaded*}
%\begin{exmp}
%If a random variable $X$ is independent of the event space $\sG$, then $\E[X\mid \sG] = \E[X]$. \\
%We may prove it by considering any set $A \in \sG$. Since, $X$ is independent of the event space $\sG$, $X$ is independent of $\Ind{A$. Then,
%\EQ{
%\E[X\Ind{A}] = \E[X]\E[\Ind{A}] = \E[\E[X]\Ind{A}].
%}
%\end{exmp}
%\end{shaded*}
%\begin{shaded*}
%\begin{exmp}
%\end{exmp}
%\end{shaded*}
\end{document}
